% 01 Cover
\begin{frame}[plain]
\begin{tikzpicture}[remember picture,overlay]
\begin{scope}[shift={(current page.north west)},x=1mm,y=-1mm]
\R{0}{0}{320}{180}{ink}
\T[white]{19}{17}{260}{16}{第七周汇报}
\T[white]{19}{49}{215}{39}{\textbf{本地长上下文Agent\\部署与推理优化}}
\T[white]{20}{111}{205}{20}{个人与公司方案、A100实现、AI Infra Book}
\T[white]{20}{157}{190}{14}{2026年10月15日}
\R{247}{48}{2}{21}{teal}\T[white]{258}{49}{57}{18}{Spark\\个人部署}
\R{247}{79}{2}{21}{blue}\T[white]{258}{80}{57}{18}{DeepSeek\\公司服务}
\R{247}{110}{2}{21}{gold}\T[white]{258}{111}{57}{18}{A100与\\推理优化}
\end{scope}\end{tikzpicture}
\end{frame}

% 02 Agenda / decisions
\begin{frame}[plain]
\StartSlide{本次汇报的四个主题}{内容概览}
\T[teal]{19}{53}{24}{29}{01}\T{54}{54}{245}{21}{\textbf{Spark个人方案}}\T[muted]{54}{67}{245}{17}{每人能用多大的上下文，生成多快，投入多少}
\Sep{54}{83}{302}
\T[blue]{19}{89}{24}{29}{02}\T{54}{90}{245}{21}{\textbf{DeepSeek公司方案}}\T[muted]{54}{103}{245}{17}{20个长会话的缓存容量、并发目标和服务器预算}
\Sep{54}{119}{302}
\T[gold]{19}{125}{24}{29}{03}\T{54}{126}{116}{21}{\textbf{A100项目}}\T[muted]{54}{140}{116}{16}{实现特点与modified vLLM对照}
\T[ink]{184}{125}{24}{29}{04}\T{218}{126}{90}{21}{\textbf{AI Infra Book}}\T[muted]{218}{140}{90}{16}{计算、缓存与任务成本}
\F{基础上下文：131K／262K。长上下文：512K／1M。总量均包含输入与输出。}
\EndSlide
\end{frame}

% 03 Personal architecture
\begin{frame}[plain]
\StartSlide{Spark运行模型，员工电脑执行工具}{个人方案}
\T{18}{48}{282}{20}{Qwen3.8-27B提供模型判断，OpenCode管理文件、工具调用与产物}
\T[teal]{23}{69}{85}{15}{员工电脑}
\T[teal]{155}{69}{72}{15}{DGX Spark 128GB}
\N{20}{81}{86}{44}{19}{\textbf{OpenCode}\\文件、Python、Excel\\员工身份与NAS权限}
\N{151}{81}{75}{44}{19}{\textbf{Qwen27模型服务}\\vLLM API}
\draw[arr] (106,92)--(151,92);\T[muted]{111}{82}{37}{12}{输入与工具结果}
\draw[back] (151,113)--(106,113);\T[muted]{111}{117}{38}{12}{答案或工具调用}
\draw[linegray,line width=.4mm] (238,70)--(238,145);
\T{250}{79}{57}{19}{\textbf{交付产物}}
\T{250}{95}{57}{17}{代码与测试\\Excel与图表\\有出处的说明}
\T[muted]{20}{138}{209}{17}{模型可共享，每个人的文件目录、账号和会话保持独立}
\F{性能数据采用FP8权重、FP8 KV Cache和MTP3。BF16 KV容量见附录B，完整配置与说明见汇报讲稿。}
\EndSlide
\end{frame}

% 04 Spark full-context estimates
\begin{frame}[plain]
\StartSlide{Spark单人生成速度与首次等待}{个人方案　性能估算}
\T[muted]{18}{47}{282}{17}{单请求接近用满上下文，预留8K输出。下表为FP8 KV＋MTP3估算}
\T[muted]{20}{66}{42}{16}{总上下文}\T[muted]{83}{66}{57}{16}{生成速度／token/s}\T[muted]{159}{66}{57}{16}{首次prefill／分钟}
\Sep{20}{77}{218}
\T{20}{85}{42}{21}{131K}\T{89}{85}{47}{24}{\V{SparkD131}}\T{169}{85}{44}{24}{\V{SparkP131}}
\T{20}{105}{42}{21}{262K}\T{89}{105}{47}{24}{\V{SparkD262}}\T{169}{105}{44}{24}{\V{SparkP262}}
\R{17}{123}{202}{19}{teal!7}\T[teal]{20}{126}{42}{21}{512K}\T[teal]{89}{126}{47}{24}{\V{SparkD512}}\T[teal]{169}{126}{44}{24}{\V{SparkP512}}
\T{20}{146}{42}{20}{1M*}\T{89}{146}{47}{22}{\V{SparkD1M}}\T{169}{146}{44}{22}{\V{SparkP1M}}
\draw[linegray,line width=.4mm] (231,65)--(231,148);
\T[teal]{244}{69}{61}{18}{\textbf{单人目标}}
\T[teal]{244}{84}{61}{31}{\textbf{20+}}
\T[muted]{244}{103}{61}{17}{token/s}
\T{244}{124}{62}{18}{512K约15\\低于当前目标}
\F{*1M速度为外推。依据：单台Spark外部实测及四档估算；实测100K约20.1、500K约14.8 token/s，详见附录A。}
\EndSlide
\end{frame}

% 05 Prefix cache example
\begin{frame}[plain]
\StartSlide{缓存命中后，只对新增输入做prefill}{个人方案　多轮请求}
\T{18}{47}{285}{19}{两种请求总上下文均接近512K，预留8K输出，本轮工具结果已包含在总量中}
\T[muted]{19}{65}{270}{15}{以下假设原输入的KV Cache和线性注意力状态已在GPU上，可直接复用}
\N[blue!6]{20}{87}{117}{24}{18}{已有输入的缓存}
\N[gold!8]{157}{87}{67}{24}{18}{新增1K工具结果}
\draw[arr] (137,99)--(157,99);\draw[arr] (224,99)--(244,99);
\T[teal]{253}{87}{53}{28}{\textbf{约\V{SparkNew1512}秒}}
\T[muted]{253}{104}{53}{13}{新增输入prefill}
\N[blue!6]{20}{125}{117}{24}{18}{已有输入的缓存}
\N[gold!8]{157}{125}{67}{24}{18}{新增8K工具结果}
\draw[arr] (137,137)--(157,137);\draw[arr] (224,137)--(244,137);
\T[teal]{253}{125}{53}{28}{\textbf{约\V{SparkNew8512}秒}}
\T[muted]{253}{142}{53}{13}{新增输入prefill}
\F{5秒与40秒是新增prefill的计算估算。首token等待还包括排队和第一步decode。缓存命中省去重复prefill，decode仍需读取KV。}
\EndSlide
\end{frame}

% 06 Pool and costs
\begin{frame}[plain]
\StartSlide{三人共享1.81M token的KV Cache}{个人方案　容量与费用}
\T{18}{48}{180}{18}{FP8 KV Cache总池约\V{Pool}M token，按请求分配}
\T[muted]{18}{67}{180}{16}{三人均分}
\R{18}{79}{58}{17}{teal!18}\R{76}{79}{58}{17}{blue!14}\R{134}{79}{58}{17}{gold!20}
\T{26}{83}{47}{17}{约600K}\T{84}{83}{47}{17}{约600K}\T{142}{83}{47}{17}{约600K}
\T[muted]{18}{107}{180}{16}{一人使用1M，另外两人共用约810K}
\R{18}{120}{96}{17}{teal!18}\R{114}{120}{39}{17}{blue!14}\R{153}{120}{39}{17}{gold!20}
\T{53}{124}{50}{17}{1M}\T{118}{124}{35}{16}{约400K}\T{157}{124}{35}{16}{约400K}
\draw[linegray,line width=.4mm] (210,47)--(210,149);
\T[teal]{226}{48}{80}{19}{\textbf{人民币预算}}
\T{226}{69}{80}{18}{硬件：5.8-6.2万}
\T{226}{86}{80}{18}{个人实施：1-2万}
\T[teal]{226}{111}{80}{27}{\textbf{6.8-8.2万}}
\T[muted]{226}{130}{80}{15}{个人硬件与实施合计}
\F{留10\%缓存余量时，三人可各配置512K。三人共享的硬件与实施合计7.8-9.9万元。价格为10月12日估算，构成见附录D。}
\EndSlide
\end{frame}

% 07 RAG concise recap
\begin{frame}[plain]
\StartSlide[blue]{公司知识库提供按权限检索的资料}{公司知识库　一页回顾}
\T{18}{48}{285}{19}{员工问答和通用Agent共用一套NAS检索接口}
\N[blue!5]{18}{80}{57}{32}{18}{\textbf{NAS}\\文件与访问权限}
\N[blue!5]{92}{80}{58}{32}{18}{\textbf{入库与索引}\\解析、向量、关键词}
\N[blue!5]{167}{80}{59}{32}{18}{\textbf{权限过滤}\\按员工身份检索}
\N[blue!5]{243}{80}{59}{32}{18}{\textbf{引用片段}\\问答或Agent使用}
\draw[arr] (75,96)--(92,96);\draw[arr] (150,96)--(167,96);\draw[arr] (226,96)--(243,96);
\T[blue]{19}{127}{135}{20}{\textbf{Onyx Enterprise}}
\T[muted]{19}{141}{135}{16}{平台、NAS适配与交付支持}
\T[blue]{170}{127}{135}{20}{\textbf{Milvus＋自建RAG}}
\T[muted]{170}{141}{135}{16}{团队维护权限同步、索引与管理接口}
\F{Word、PDF、PPT、Excel共用入库流程。内容进入重排和模型前完成权限过滤，原文入口复核权限。详见公司知识库方案。}
\EndSlide
\end{frame}

% 08 Enterprise architecture
\begin{frame}[plain]
\StartSlide[blue]{公司Agent共享模型，各自执行任务}{公司方案　DeepSeek-V4.1-Flash}
\N[blue!5]{18}{74}{70}{39}{19}{\textbf{Hermes会话}\\员工／项目独立\\Skills与任务记录}
\N[blue!5]{125}{74}{80}{39}{19}{\textbf{共享模型服务}\\DeepSeek-V4.1-Flash}
\draw[arr] (88,85)--(125,85);\draw[back] (125,103)--(88,103);
\T[muted]{91}{70}{33}{12}{输入／工具结果}\T[muted]{92}{115}{35}{12}{答案／工具调用}
\N{241}{59}{63}{30}{18}{GPU：权重与\\各会话KV Cache}
\N{241}{103}{63}{30}{18}{主机内存：\\约203GB Engram}
\draw[blue,line width=.5mm] (205,94)--(223,94);
\draw[blue,line width=.5mm] (241,74)--(223,74)--(223,118)--(241,118);
\T{19}{132}{190}{18}{工具在员工电脑或CPU服务器执行\\MCP复用公司知识库和业务接口}
\T[blue]{19}{47}{285}{20}{目标：20个同时活跃的任务，每个最多1M上下文}
\F{Engram表按模型实例计。KV Cache按所有保留会话计，包括正在等待工具的会话。用户身份随检索请求传递。}
\EndSlide
\end{frame}

% 09 Capacity and instance alternatives
\begin{frame}[plain]
\StartSlide[blue]{20个会话分别能分到多少上下文}{公司方案　缓存容量估算}
\T[muted]{18}{46}{287}{16}{按公开启动池及图捕获后预算估算，再预留10\%缓存容量}
\T[muted]{19}{66}{105}{16}{配置}\T[muted]{129}{66}{66}{16}{每会话上下文}
\Sep{19}{78}{197}
\T{19}{86}{104}{20}{8×H100，TP8}\T{137}{86}{57}{24}{约\V{H100K}K}
\T{19}{108}{104}{20}{4×H200，TP4}\T{137}{108}{57}{24}{约\V{H200fourK}K}
\R{16}{129}{181}{20}{blue!6}\T[blue]{19}{133}{104}{20}{8×H200}\T[blue]{137}{132}{57}{24}{1M}
\draw[linegray,line width=.4mm] (210,63)--(210,149);
\T[blue]{224}{66}{82}{19}{\textbf{同机两种部署}}
\T{224}{85}{82}{18}{TP8：一个实例\\一组队列，20人}
\T[muted]{224}{108}{82}{14}{Engram主存约203GB}
\T{224}{124}{82}{18}{双TP4：两组各10人}
\T[muted]{224}{140}{82}{14}{两份Engram约406GB}
\F{1M总量含8K输出预留。H100的约242K低于20人各262K所需容量。该页描述缓存容量，20路满1M的速度另需对应负载测量。}
\EndSlide
\end{frame}

% 10 Resident vs active
\begin{frame}[plain]
\StartSlide[blue]{等待工具的请求也占用KV Cache}{公司方案　并发与状态}
\T[blue]{21}{55}{126}{30}{\textbf{20个请求在decode}}
\T{21}{77}{126}{18}{使用GPU计算与KV Cache}
\T[gold]{175}{55}{126}{30}{\textbf{20个请求在等工具}}
\T{175}{77}{126}{18}{暂不生成，继续保留KV Cache}
\draw[linegray,line width=.5mm] (160,52)--(160,98);
\Sep{20}{101}{302}
\T{21}{111}{145}{24}{\textbf{缓存按40份上下文计算}}
\T[muted]{21}{132}{145}{17}{闲置会话可移至CPU内存\\继续任务时，再复制回GPU}
\T[blue]{191}{109}{115}{24}{\textbf{生成速度目标}}
\T{191}{131}{115}{18}{20路 × 20 token/s\\= 400 token/s系统输出}
\F{若每路10 token/s，系统目标为200。当前TP4实现恢复一个1M会话的四份全局KV约需传输7.3GB，另计复制等待。}
\EndSlide
\end{frame}

% 11 Enterprise costs and software
\begin{frame}[plain]
\StartSlide[blue]{服务器预算与公司推理引擎}{公司方案　人民币万元}
\T[muted]{19}{49}{123}{16}{完整GPU服务器}\T[muted]{160}{49}{62}{16}{交付预算估算}\T[muted]{242}{49}{63}{16}{主机内存}
\Sep{19}{61}{302}
\T{19}{68}{132}{20}{8×A100 80GB SXM，翻新}\T{166}{68}{60}{23}{\V{A100Price}}\T{250}{70}{54}{18}{1TB}
\T{19}{91}{132}{20}{8×H100 80GB SXM}\T{166}{91}{60}{23}{\V{H100Price}}\T{250}{93}{54}{18}{1.5TB}
\T{19}{114}{132}{20}{4×H200 141GB NVL}\T{166}{114}{60}{23}{\V{H200fourPrice}}\T{250}{116}{54}{18}{1TB}
\R{16}{134}{290}{19}{blue!6}\T[blue]{19}{137}{132}{20}{8×H200 141GB SXM}\T[blue]{166}{137}{60}{23}{\V{H200eightPrice}}\T[blue]{250}{139}{54}{18}{1.5TB}
\F{10月12日价格与估算，USD/CNY=7.00。集中工具执行另加12-26万元。公司采用A100时选择固定版本modified vLLM，详见A100部分。}
\EndSlide
\end{frame}

% 12 A100 intro and actual topology
\begin{frame}[plain]
\StartSlide[gold]{A100项目：新模型在四卡机器上的移植}{shi3z工程案例}
\T{18}{46}{286}{19}{稠密部分按层指定owner GPU，路由专家分散到四张卡上}
\N[gold!7]{19}{79}{77}{32}{18}{\textbf{当前层owner}\\attention与路由}
\N[gold!7]{129}{79}{77}{32}{18}{\textbf{四卡专家计算}\\按路由执行并汇总}
\draw[arr] (96,89)--(129,89);\draw[back] (129,102)--(96,102);
\N{239}{79}{64}{32}{18}{\textbf{下一层}\\必要时切换owner}
\draw[arr] (58,79)--(58,67)--(270,67)--(270,79);
\T[muted]{112}{60}{112}{12}{本层结果进入下一层}
\draw[arr] (270,111)--(270,122)--(58,122)--(58,111);
\T[muted]{107}{125}{134}{14}{继续下一层，直到40层完成}
\T[gold]{19}{142}{285}{18}{\textbf{CPU内存放Engram，GPU保存其余权重和会话缓存}}
\F{项目主配置：4×A100 80GB PCIe，两组NVLink卡对，组间PCIe。源码固定580568e。箭头表示一层内协作与层间推进。}
\EndSlide
\end{frame}

% 13 Contribution attribution
\begin{frame}[plain]
\StartSlide[gold]{模型能力、通用优化与项目实现}{A100案例　贡献归属}
\T[muted]{20}{49}{70}{16}{来源}\T[muted]{108}{49}{195}{16}{具体内容}
\Sep{20}{61}{302}
\T[gold]{20}{70}{75}{22}{\textbf{DeepSeek模型}}
\T{108}{69}{195}{19}{原生FP4／FP8权重、CED、共享KV、DSpark\\决定计算路径与原生存储格式}
\Sep{20}{97}{302}
\T[gold]{20}{106}{75}{22}{\textbf{成熟引擎}}
\T{108}{105}{195}{19}{并行、CUDA Graph、算子融合、前缀缓存\\提供已有的执行与服务方法}
\Sep{20}{128}{302}
\T[gold]{20}{136}{75}{22}{\textbf{项目实现}}
\T{108}{135}{195}{18}{A100算子适配、按NVLink卡对通信\\保存与载入会话缓存}
\F{学习重点：硬件适配、通信安排和缓存管理。来源：DeepSeek模型卡、vLLM文档、shi3z源码及modified vLLM测试记录。}
\EndSlide
\end{frame}

% 14 Public performance comparison
\begin{frame}[plain]
\StartSlide[gold]{公开测速保留各自的硬件与负载}{A100案例　与modified vLLM对照}
\T[muted]{18}{47}{283}{16}{短输入结果，单位token/s。四卡PCIe与八卡SXM分别列示}
\T[muted]{19}{64}{84}{16}{实现与硬件}\T[muted]{115}{64}{82}{16}{单路结果}\T[muted]{220}{64}{84}{16}{多路结果}
\Sep{19}{76}{302}
\T{19}{83}{85}{18}{shi3z\\4×A100 PCIe}
\T{115}{83}{90}{18}{代码HTTP：50-53\\MTP实验：82-93}
\T{220}{83}{84}{18}{64路端到端：167.5\\引擎峰值：475.6}
\Sep{19}{109}{302}
\T{19}{116}{85}{18}{modified vLLM\\8×A100 SXM}
\T{115}{116}{90}{18}{DSpark5 decode\\110-140}
\T{220}{116}{84}{18}{32并发总输出\\1,040-1,379}
\Sep{19}{142}{302}
\T[gold]{19}{146}{282}{17}{公司A100采用modified vLLM，复用调度、缓存、接口与监控}
\F{第一行来自作者不同实验。第二行为keplerzip历史TP8、短输入1,024输出的操作者回传结果。缺少两套引擎的同机同负载直接对照。}
\EndSlide
\end{frame}

% 15 Controlled ablation
\begin{frame}[plain]
\StartSlide[gold]{同机消融量出每项改动的额外收益}{modified vLLM　八卡A800}
\T[muted]{18}{47}{286}{16}{八卡A800，128并发，每档两次；总输出token/s。基线含DSpark、Marlin和CUDA Graph}
\T{19}{69}{67}{17}{Marlin dense\\NCCL EP8}
\R{91}{73}{156.00}{12}{linegray}\T{254}{71}{52}{23}{\V{AblationBase}}
\T{19}{99}{67}{17}{大batch稠密层\\改走BF16算子}
\R{91}{103}{160.59}{12}{gold!50}\T{254}{101}{52}{23}{\V{AblationDense}}
\T[gold]{91}{119}{110}{15}{相对上一档 +2.94\%}
\T{19}{130}{67}{17}{再加入\\定制EP8通信}
\R{91}{134}{163.03}{12}{gold}\T{260}{131}{48}{23}{\V{AblationComm}}
\T[gold]{91}{150}{180}{13}{再增加 +1.52\%，两项合计约4.5\%}
\F{条形从零起算，长度与输出吞吐成比例。输入78-90、输出1,024 tokens。来源：Tokha233公开消融与测试协议。}
\EndSlide
\end{frame}

% 16 Correct native memory budget
\begin{frame}[plain]
\StartSlide[gold]{四卡显存从官方510GB原生制品起算}{A100案例　权重与运行空间}
\T{18}{47}{287}{20}{官方checkpoint约\V{NativeGB}GB，原生采用FP4专家与FP8稠密权重}
\R{19}{77}{110}{28}{blue!16}\R{129}{77}{167}{28}{teal!17}
\T[blue]{27}{82}{95}{22}{Engram \V{EngramGB}GB}
\T[teal]{143}{82}{149}{22}{其余权重 \V{GPUWeightsGB}GB}
\T[muted]{27}{109}{97}{17}{放在主机内存}
\T[muted]{143}{109}{149}{17}{约\V{GPUWeightsGiB}GiB，分配到GPU}
\T[gold]{19}{135}{119}{28}{\textbf{四卡平均\V{PerCard}GiB}}
\T{166}{132}{137}{18}{每卡还需KV Cache、草稿与工作区\\具体分片决定最紧张的一张卡}
\F{原生权重与Engram卸载是模型和框架已有路径。项目进一步处理算子、分配和临时张量。依据：官方tensor索引与机制复算。}
\EndSlide
\end{frame}

% 17 Book critical path
\begin{frame}[plain]
\StartSlide{一次Agent任务的完整耗时}{AI Infra Book　第8章与第11章}
\T{18}{47}{286}{20}{以财报任务为例，用户等待的是输入处理、生成和工具执行的总时间}
\N{18}{79}{49}{30}{18}{请求排队}
\N[teal!7]{82}{79}{54}{30}{18}{输入prefill}
\N[blue!7]{151}{79}{54}{30}{18}{生成工具调用}
\N[gold!8]{220}{79}{82}{30}{18}{Python／Excel等工具}
\draw[arr] (67,94)--(82,94);\draw[arr] (136,94)--(151,94);\draw[arr] (205,94)--(220,94);
\draw[back] (261,109)--(261,123)--(109,123)--(109,109);
\T[muted]{127}{127}{119}{14}{工具结果返回，进入下一轮}
\T{19}{143}{284}{20}{读取100GB，若有效带宽为1TB/s，仅传输约需0.1秒}
\F{每阶段分别估算计算量／有效算力、字节数／有效带宽，再计入前后依赖和排队。来源：AI Infra Book第8、9、11章。}
\EndSlide
\end{frame}

% 18 Book batching + context
\begin{frame}[plain]
\StartSlide{合批分摊权重读取，KV按请求分别读取}{AI Infra Book　带宽与上下文}
\T[muted]{18}{47}{286}{17}{稠密模型的简化读取量，按每个输出token计}
\T[teal]{24}{67}{279}{25}{\textbf{权重读取量 ÷ batch内请求数 ＋ 该请求的KV读取量}}
\Sep{20}{94}{302}
\T{20}{107}{120}{20}{上下文较短\\权重读取占比大，合批收益明显}
\T{175}{107}{126}{20}{上下文较长\\每个请求要读取更多KV}
\T[teal]{20}{140}{281}{20}{Qwen从512K增至1M：KV约1.9倍，全注意力冷prefill计算约3.7倍}
\F{1.9倍指KV容量，3.7倍指全注意力部分的计算量。缓存命中可省重复prefill，新增输入仍需读取已有KV。来源：Book算例。}
\EndSlide
\end{frame}

% 19 Book tool residence
\begin{frame}[plain]
\StartSlide{CPU不忙时，任务环境仍可能占用大量内存}{AI Infra Book　资源调度算例}
\T[muted]{18}{46}{286}{16}{每10秒到达1项任务，每项3轮。每轮模型60秒、工具10秒，工具实际消耗4核秒}
\T[muted]{19}{65}{121}{16}{资源}\T[muted]{149}{65}{99}{16}{计算}\T[muted]{259}{65}{45}{16}{平均需求}
\Sep{19}{77}{302}
\T{19}{84}{122}{19}{尚未完成的任务}\T{149}{84}{99}{18}{210秒 ÷ 10秒}\T[teal]{259}{83}{45}{23}{21项}
\T{19}{106}{122}{19}{处于模型阶段的任务}\T{149}{106}{99}{18}{180秒 ÷ 10秒}\T[teal]{259}{105}{45}{23}{18项}
\T{19}{128}{122}{19}{工具CPU}\T{149}{128}{99}{18}{12核秒 ÷ 10秒}\T[teal]{259}{127}{45}{23}{1.2核}
\T{19}{146}{122}{18}{工具环境内存}\T{149}{146}{99}{18}{21项 × 2GiB}\T[teal]{259}{145}{45}{23}{42GiB}
\F{教学假设：稳定到达、忽略排队，每个环境占2GiB并保留至任务结束。模型KV、工具CPU与环境内存分别预算。}
\EndSlide
\end{frame}

% 20 Summary
\begin{frame}[plain]
\StartSlide{配置选择与使用方式}{汇报总结}
\T[teal]{19}{52}{73}{22}{\textbf{个人Spark}}
\T{105}{51}{198}{20}{Qwen27，工具在员工电脑执行\\512K估算约15 token/s，共享KV约1.81M tokens}
\Sep{19}{82}{302}
\T[blue]{19}{94}{73}{22}{\textbf{公司DeepSeek}}
\T{105}{93}{198}{20}{八卡H200为主要候选，比较TP8与双TP4\\缓存可容纳20个1M会话，按活跃及等待任务计容量}
\Sep{19}{124}{302}
\T[gold]{19}{136}{73}{22}{\textbf{公司A100}}
\T{105}{135}{198}{20}{采用固定版本modified vLLM\\学习shi3z的实现思路，按同机测试选择补丁}
\F{个人20+ token/s目标仍高于512K的当前估算。公司20路满1M速度需匹配负载验证。最终经济性按合格任务的时间与成本比较。}
\EndSlide
\end{frame}

% 21 Appendix measured single Spark
\begin{frame}[plain]
\StartSlide{Spark单请求实测与四档估算的关系}{附录A　测试口径}
\T[muted]{18}{47}{286}{16}{单GB10，FP8权重＋FP8 KV，MTP3，随机token输入，前缀缓存关闭}
\T[muted]{19}{65}{65}{15}{实际输入token}\T[muted]{100}{65}{48}{15}{输出token}\T[muted]{169}{65}{54}{15}{平均TTFT／秒}\T[muted]{245}{65}{60}{15}{decode／token/s}
\Sep{19}{76}{302}
\T{19}{83}{65}{18}{1,024}\T{100}{83}{48}{18}{2,048}\T{176}{83}{47}{18}{0.6}\T{254}{83}{50}{18}{20.53}
\T{19}{97}{65}{18}{50,000}\T{100}{97}{48}{18}{2,048}\T{176}{97}{47}{18}{63.1}\T{254}{97}{50}{18}{21.01}
\T{19}{111}{65}{18}{100,000}\T{100}{111}{48}{18}{2,048}\T{176}{111}{47}{18}{144.8}\T{254}{111}{50}{18}{20.11}
\T{19}{125}{65}{18}{300,000}\T{100}{125}{48}{18}{2,048}\T{176}{125}{47}{18}{658.9}\T{254}{125}{50}{18}{17.23}
\T{19}{139}{65}{18}{500,000}\T{100}{139}{48}{18}{2,048}\T{176}{139}{47}{18}{1,470.1}\T{254}{139}{50}{18}{14.82}
\F{另有999,936输入＋64输出的1M容量测试：TTFT4,815.8秒，短段decode9.83。四档估算预留8K输出，与本表实际长度分别记录。}
\EndSlide
\end{frame}

% 22 Appendix cache budget
\begin{frame}[plain]
\StartSlide{缓存池、运行缓冲与上下文容量}{附录B　容量计算}
\T[muted]{19}{49}{87}{15}{公开配置}\T[muted]{113}{49}{58}{15}{已分配KV／卡}\T[muted]{185}{49}{55}{15}{实例token池}\T[muted]{251}{49}{52}{15}{严格预算／卡}
\Sep{19}{61}{302}
\T{19}{70}{88}{18}{H100 TP8}\T{115}{70}{58}{18}{12.36GiB}\T{188}{70}{54}{18}{6.435M}\T{253}{70}{51}{18}{10.34GiB}
\T{19}{92}{88}{18}{H200 TP4}\T{115}{92}{58}{18}{35.16GiB}\T{188}{92}{54}{18}{18.327M}\T{253}{92}{51}{18}{32.49GiB}
\T{19}{114}{88}{18}{H200 TP8}\T{115}{114}{58}{18}{62.86GiB}\T{188}{114}{54}{18}{29.498M}\T{253}{114}{51}{18}{60.20GiB}
\Sep{19}{135}{302}
\T{19}{143}{286}{17}{Spark FP8池：\V{PoolExact} tokens。BF16容量约减半，可容纳单人512K或三人各262K}
\F{公司主表对graph后建议预算线性调整，再留10\%tokens。Spark BF16为同池容量近似。预留空闲tokens不会释放已分配的物理内存。}
\EndSlide
\end{frame}

% 23 Appendix public enterprise performance
\begin{frame}[plain]
\StartSlide{公司服务的公开多轮回放参照}{附录C　AgentX性能}
\T[muted]{18}{46}{286}{16}{平均输入140K-157K，前缀命中约95\%，预设DSpark平均接受长度3.51}
\T[muted]{19}{65}{94}{15}{配置与回放会话数}\T[muted]{128}{65}{72}{15}{总输出／token/s}\T[muted]{222}{65}{82}{15}{TTFT均值／p95 秒}
\Sep{19}{77}{302}
\T{19}{85}{99}{18}{H100 TP8，20}\T{143}{85}{65}{20}{718.48}\T{231}{85}{73}{18}{2.05／6.98}
\T{19}{100}{99}{18}{H200 TP4，16}\T{143}{100}{65}{20}{614.76}\T{231}{100}{73}{18}{0.83／3.03}
\T{19}{115}{99}{18}{H200 TP4，32}\T{143}{115}{65}{20}{922.61}\T{231}{115}{73}{18}{1.47／5.12}
\T{19}{130}{99}{18}{H200 TP8，16}\T{143}{130}{65}{20}{661.15}\T{231}{130}{73}{18}{0.65／2.14}
\T{19}{145}{99}{18}{H200 TP8，32}\T{143}{145}{65}{20}{1,052.18}\T{231}{145}{73}{18}{0.85／2.63}
\F{端到端总输出含输入处理与调度。回放包含生成与等待工具两个阶段；该表用于140K-157K输入下的性能参照。来源：vLLM AgentX记录。}
\EndSlide
\end{frame}

% 24 Appendix price scope
\begin{frame}[plain]
\StartSlide{预算包含什么，哪些费用单独计算}{附录D　价格与交付}
\T[blue]{20}{51}{130}{23}{\textbf{Spark预算}}
\T{20}{70}{132}{18}{美国页面价：6,950美元\\预算汇率：1美元＝7.00元\\税费、运输与支持准备另加\\国内交付估算：5.8-6.2万元}
\T[muted]{20}{119}{132}{16}{个人实施：4-8人天 × 2,500元\\三人实施：8-15人天 × 2,500元}
\draw[linegray,line width=.4mm] (161,49)--(161,148);
\T[blue]{178}{51}{128}{23}{\textbf{GPU服务器预算}}
\T{178}{70}{127}{18}{完整主机挂牌或组件BOM\\加13\%税费准备、10\%价格预留\\交付及本地支持单列\\CPU、RAM、电源按已含配置计一次}
\T[muted]{178}{119}{126}{16}{工具CPU节点另加12-26万元\\机房改造、软件授权与长期运维另计}
\F{价格观察／估算日期为2026年10月12日。人民币范围为预算，不是境内成交报价。准确SKU、供应商交付路线与许可文件随采购核对。}
\EndSlide
\end{frame}

% 25 Reading / references
\begin{frame}[plain]
\StartSlide{主要来源与继续阅读}{附录E}
\T[teal]{20}{51}{135}{22}{\textbf{模型与推理引擎}}
\T{20}{72}{135}{18}{\href{https://huggingface.co/deepseek-ai/DeepSeek-V4.1-Flash}{DeepSeek-V4.1-Flash模型与权重}\\[5pt]\href{https://github.com/vllm-project/vllm}{vLLM源码与Engram实现}\\[5pt]\href{https://github.com/shi3z/deepseekv4.1-A100-custom}{shi3z A100定制runtime}\\[5pt]\href{https://github.com/Tokha233/deepseek-v4.1-flash-a100-turbo}{modified vLLM的A800测试与消融}}
\T[blue]{178}{51}{126}{22}{\textbf{AI Infra Book}}
\T{178}{72}{126}{18}{\href{https://github.com/bojieli/ai-infra-book}{李博杰：深入理解AI Infra}\\[7pt]第8章：推理优化\\[5pt]第9章：分布式推理\\[5pt]第11章：资源调度与运行环境}
\Sep{20}{134}{302}
\T[muted]{20}{143}{281}{16}{完整解释见本次汇报讲稿与场景方案，知识库保留原文、版本和计算输入}
\F{资料截至2026年10月14日。Spark数据来自外部设备测试；AI Infra Book与A100源码采用本次研究归档版本。}
\EndSlide
\end{frame}
